\documentclass[11pt]{article}

\usepackage{amsmath,amssymb,amsthm,mathtools}
\usepackage[margin=1in]{geometry}
\usepackage{url}

\newtheorem{theorem}{Theorem}[section]
\newtheorem{lemma}[theorem]{Lemma}
\newtheorem{corollary}[theorem]{Corollary}
\newtheorem{remark}[theorem]{Remark}

\newcommand{\dd}{\,d}
\newcommand{\Gap}{\Gamma_H}

\begin{document}

\section{Purpose and the perturbation families}

For a finite graph $H$, put
\[
 \Gamma_H(W)=t(H,W)-\Phi_H(p(W)),
 \qquad
 \Phi_H(p)=(1-p)^{v(H)}
 \chi_H\!\left(\frac1{1-p}\right).
\]
This note combines the two tests requested in the accompanying research:
arbitrary changes around a complete multipartite graphon and the
high-density limit $p\uparrow1$.  Its main theorem controls every arbitrary
measurable complement graphon of constant degree $q$ when $q$ is
sufficiently small.  It gives a no-hit theorem, not a proof of the
conjecture for all graphons.

Let $V=1-W$ be the complement.  We first assume only
\begin{equation}
 0\leq V\leq1,
 \qquad d_V(x)=\int V(x,y)\dd y=q
 \quad\hbox{for almost every }x.
 \label{eq:regular}
\end{equation}
Thus $V$ may be non-step and the underlying probability space arbitrary.
The balanced Turan complements are singular equality cases of this class.

Fix an integer $r\geq2$ and split the probability space into $r$ sets of
measure $1/r$.  Let $B=(B_{ij})$ be an arbitrary symmetric doubly stochastic
$r\times r$ matrix, and define the complement graphon $V_B$ and graphon
$W_B$ by
\[
 V_B(x,y)=B_{ij},\qquad W_B=1-V_B
 \quad (x\in\Omega_i,\ y\in\Omega_j).
\]
Every row sum of $B$ is one, so $V_B$ is a special case of
\eqref{eq:regular}, of degree
$q=1/r$ and $p(W_B)=1-1/r$.  The choice $B=I_r$ is the balanced Turan
graphon complement.  The choice $B=J_r/r$ makes $W_B$ constant.  More
generally this family contains every feasible degree-preserving coarse
perturbation of $I_r$, and the whole segment
$(1-\theta)I_r+\theta J_r/r$, $0\leq\theta\leq1$.  Thus the test is neither
restricted to one infinitesimal direction nor to multipartite graphons.

\section{Exact graphic-rank expansion}

For $A\subseteq E(H)$, write $H_A=(V(H),A)$ and let
$r_H(A)=v(H)-c(H_A)$ be its graphic-matroid rank.  Inclusion--exclusion and
Whitney's expansion of the chromatic polynomial give
\begin{equation}
 \Gamma_H(1-V)=
 \sum_{A\subseteq E(H)}(-1)^{|A|}
 \bigl(t(H_A,V)-q^{r_H(A)}\bigr)
 \label{eq:rank-expansion}
\end{equation}
whenever $\int V=q$.

\begin{lemma}[Forests are exact]\label{lem:forest}
If $V$ is regular of degree $q$, then
$t(F,V)=q^{e(F)}$ for every forest $F$.  If $C$ is a connected graph on
$v$ vertices, then
\[
 0\leq q^{v-1}-t(C,V).
\]
\end{lemma}

\begin{proof}
Repeatedly integrate a leaf to obtain a factor $q$; isolated vertices give
factor one.  For the second assertion, choose a spanning tree $T$ of $C$.
Since $0\leq V\leq1$, deleting the edges outside $T$ can only increase the
density, so $t(C,V)\leq t(T,V)=q^{v-1}$.
\end{proof}

For any regular $V$, define the exact transitivity defect
\begin{equation}
 D_V:=t(P_3,V)-t(K_3,V)=q^2-t(K_3,V)\geq0.
 \label{eq:D}
\end{equation}
Nonnegativity follows pointwise from
\[
 D_V=\int V(x,y)V(y,z)(1-V(x,z))
       \dd x\dd y\dd z.
\]
For the stochastic-block specialization,
\begin{equation}
 D_{V_B}=\frac{r-\operatorname{tr}(B^3)}{r^3}.
 \label{eq:block-D}
\end{equation}

\section{A spectral bound for every cyclic component}

\begin{lemma}[Cycle defect]\label{lem:cycle}
For every integer $\ell\geq3$,
\begin{equation}
 0\leq q^{\ell-1}-t(C_\ell,V)
 \leq (\ell-2)q^{\ell-3}D_V.
 \label{eq:cycle-defect}
\end{equation}
\end{lemma}

\begin{proof}
If $q=0$, then $V=0$ almost everywhere and the assertion is immediate.
Suppose $q>0$.  Let $T_V$ be the self-adjoint Hilbert--Schmidt integral
operator with kernel $V$.  The Schur test and \eqref{eq:regular} give
$\lVert T_V\rVert_{2\to2}\leq q$.  The constant function is an eigenvector
of eigenvalue $q$; write the remaining eigenvalues as $q\lambda_i$, where
$\lambda_i\in[-1,1]$.  Moreover
\[
 1+\sum_i\lambda_i^2
 =\frac{\lVert T_V\rVert_{\rm HS}^2}{q^2}
 =\frac{\int V^2}{q^2}
 \leq\frac1q.
\]
Put
$E=1/q-1-\sum_i\lambda_i^2\geq0$.  The trace formula for powers of a
Hilbert--Schmidt kernel gives
\[
 \frac{q^{\ell-1}-t(C_\ell,V)}{q^\ell}
 =E+\sum_i(\lambda_i^2-\lambda_i^\ell),
\]
and
\[
 \frac{D_V}{q^3}=E+\sum_i(\lambda_i^2-\lambda_i^3).
\]
For $\lambda\in[-1,1]$,
\[
 0\leq\lambda^2-\lambda^\ell
 \leq(\ell-2)(\lambda^2-\lambda^3).
\]
After division by $\lambda^2$ when nonzero, this is the elementary
geometric-series inequality
$0\leq(1-\lambda^{\ell-2})/(1-\lambda)\leq\ell-2$; the endpoint cases
follow by continuity.  Since $\ell-2\geq1$, summing this inequality and
retaining $E$ proves \eqref{eq:cycle-defect}.  The trace sums converge
absolutely in every power at least two.
\end{proof}

For a connected graph $C$, let
$\beta(C)=e(C)-v(C)+1$ be its cyclomatic number.

\begin{lemma}[Connected cyclic defect]\label{lem:connected}
If $C$ is connected, has $v\geq3$ vertices, and $\beta(C)>0$, then
\begin{equation}
 0\leq q^{v-1}-t(C,V)
 \leq v\,\beta(C)q^{v-3}D_V.
 \label{eq:connected-defect}
\end{equation}
\end{lemma}

\begin{proof}
Choose a spanning tree $T$.  Since
\[
 1-\prod_{e\in E(C)\setminus E(T)}V_e
 \leq\sum_{e\in E(C)\setminus E(T)}(1-V_e),
\]
the defect from $t(T,V)=q^{v-1}$ is at most the sum of one term for each
of the $\beta(C)$ extra edges.  For such an edge, its endpoints are joined
in $T$ by a path with $\ell-1$ edges, where $3\leq\ell\leq v$ after the
extra edge closes the path to a cycle $C_\ell$.  Integrating every tree
branch off this path gives the factor $q^{v-\ell}$.  The remaining integral
is exactly
\[
 q^{v-\ell}\bigl(q^{\ell-1}-t(C_\ell,V)\bigr)
 \leq (\ell-2)q^{v-3}D_V
 \leq vq^{v-3}D_V
\]
by Lemma~\ref{lem:cycle}.  Sum over the extra edges.
\end{proof}

For a possibly disconnected $H_A$, define
\[
 c(A)=\sum_{C}v(C)\beta(C),
\]
where the sum is over its nontrivial cyclic components.  Tree components
make no contribution.

\begin{lemma}[Arbitrary edge-subset defect]\label{lem:subset}
If $R=r_H(A)\geq2$, then
\begin{equation}
 0\leq q^R-t(H_A,V)
 \leq c(A)q^{R-2}D_V.
 \label{eq:subset-defect}
\end{equation}
\end{lemma}

\begin{proof}
The density and the comparison monomial factor over the nontrivial
components.  Telescope the difference of these two products.  Lemma
\ref{lem:forest} makes every tree-component difference zero and bounds
every unchanged component density by its forest value.  If a cyclic
component $C$ has $v$ vertices, Lemma~\ref{lem:connected} bounds its
difference by $v\beta(C)q^{v-3}D_V$.  The comparison powers from all other
components contribute $q^{R-(v-1)}$.  Their product is
$v\beta(C)q^{R-2}D_V$.  Summing over cyclic components gives
\eqref{eq:subset-defect}.
\end{proof}

\section{Uniform regular-complement theorem}

Let $N_3(H)$ be the number of triangles of $H$ and put
\begin{equation}
 \Xi_H=\sum_{A\subseteq E(H):\ r_H(A)\geq3}c(A).
 \label{eq:Xi}
\end{equation}

\begin{theorem}[Arbitrary regular-complement stability]
\label{thm:regular}
Let $V$ be any graphon satisfying \eqref{eq:regular}, and put $W=1-V$.
Then
\begin{equation}
 \Gamma_H(W)\geq
 \bigl(N_3(H)-\Xi_Hq\bigr)D_V.
 \label{eq:main-bound}
\end{equation}
Consequently, if $N_3(H)>0$ and
\begin{equation}
 0<q<\frac{N_3(H)}{\Xi_H},
 \label{eq:q-range}
\end{equation}
then $\Gamma_H(W)\geq0$.  Equality occurs precisely when $V$ is a union of
equal diagonal clique blocks of mass $q$; necessarily $q=1/r$ for an
integer $r$, and $W=T_r$ up to null sets.  Every other regular complement
has strict inequality.
\end{theorem}

\begin{proof}
In \eqref{eq:rank-expansion}, ranks zero and one vanish by
Lemma~\ref{lem:forest}.  The two-edge forests at rank two also vanish.  Each
triangle has three edges and contributes exactly $D_V$, so the full
rank-two contribution is $N_3(H)D_V$.

For every $A$ of rank $R\geq3$, Lemma~\ref{lem:subset} bounds the absolute
value of its summand by
\[
 c(A)q^{R-2}D_V\leq c(A)qD_V,
\]
because $q\leq1$ and $R\geq3$.  Discarding every sign and summing these
bounds proves \eqref{eq:main-bound}.

Under \eqref{eq:q-range}, the coefficient in \eqref{eq:main-bound} is
positive.  Thus equality requires $D_V=0$.  The universal transitivity
identity in \eqref{eq:D} then gives, for almost every triple,
\[
 V(x,y)V(y,z)>0\quad\Longrightarrow\quad V(x,z)=1.
\]
For a typical $x$, put $N_x=\{z:V(x,z)>0\}$.  If $y\in N_x$ is typical,
apply the implication to $(z,x,y)$: $V(z,y)=1$ for almost every
$z\in N_x$.  Regularity and $0\leq V\leq1$ now give
\[
 q=d_V(x)\leq\mu(N_x)\leq d_V(y)=q.
\]
Thus the positive-support classes are clique blocks of measure $q$.
They cover the probability space up to a null set, so their number is
$r=1/q$.  Hence $V$ is the balanced clique-block complement and $W=T_r$.
The converse is immediate from the exact Turan equality.
\end{proof}

The theorem controls every exactly regular measurable complement, not just
a matrix family or a neighbourhood of a Turan graphon.  Specializing to
$V_B$ and $q=1/r$ gives
\begin{equation}
 \Gamma_H(W_B)\geq
 \left(N_3(H)-\frac{\Xi_H}{r}\right)D_{V_B}.
 \label{eq:block-bound}
\end{equation}
Thus every symmetric doubly stochastic coarse perturbation is nonnegative
when $r>\Xi_H/N_3(H)$, and equality forces $B=I_r$.

\section{The full Turan-to-constant segment}

The most canonical path inside the preceding matrix family admits a much
stronger, full-range certificate.  Put
\[
 B_\theta=(1-\theta)I_r+\frac{\theta}{r}J_r,
 \qquad 0\leq\theta\leq1.
\]
Then $B_0=I_r$, while $B_1=J_r/r$.  Thus $W_{B_0}=T_r$ and
$W_{B_1}$ is the constant graphon of density $1-1/r$.  Every intermediate
graphon has the same density.  Its two cell values are
\[
 W_{B_\theta}(x,y)=
 \begin{cases}
  \theta(1-1/r),&x,y\text{ in the same cell},\\
  1-\theta/r,&x,y\text{ in different cells}.
 \end{cases}
\]

Let $m=e(H)$ and $\beta(A)=|A|-r_H(A)$.  Expanding each edge factor according
to whether its endpoint colours agree gives the exact polynomial
\begin{align}
 G_H(r,\theta)
 &:=r^m\Gamma_H(W_{B_\theta}) \notag\\
 &=\sum_{A\subseteq E(H)}(-1)^{|A|}
 \left[
 r^{\beta(A)}(r-\theta)^{m-|A|}(1-\theta)^{|A|}
 -r^{m-r_H(A)}
 \right].
 \label{eq:interpolation-polynomial}
\end{align}
Indeed, before multiplication by $r^m$, fixing the edges in $A$ to have
equal endpoint colours leaves $r^{c(H_A)}$ colour assignments out of
$r^{v(H)}$, hence the factor $r^{-r_H(A)}$.  Equation
\eqref{eq:interpolation-polynomial} then follows.  At $\theta=0$ it vanishes
term by term, so $\theta$ divides $G_H$ in $\mathbb Z[r,\theta]$.

\begin{theorem}[Exact interpolation certificates]
\label{thm:interpolation}
For each of the $29$ currently open Atlas graphs $H$, every integer
$r\geq\chi(H)-1$, and every $\theta\in[0,1]$,
\[
 \Gamma_H(W_{B_\theta})\geq0.
\]
Equality holds at $\theta=0$.  The inequality is strict for
$0<\theta\leq1$.
\end{theorem}

\begin{proof}
Put $k=\chi(H)-1$, $y=r-k$, and
\[
 Q_H(y,\theta)=\frac{G_H(y+k,\theta)}{\theta}.
\]
If $d=\deg_\theta Q_H$, write its native-degree Bernstein expansion
\begin{equation}
 Q_H(y,\theta)=
 \sum_{i=0}^d b_{H,i}(y)\binom di
 \theta^i(1-\theta)^{d-i}.
 \label{eq:interpolation-bernstein}
\end{equation}
Exact expansion of \eqref{eq:interpolation-polynomial} shows, for $28$ of
the rows, that every $b_{H,i}(y)$ has nonnegative rational monomial
coefficients.  None is the zero polynomial.  For Atlas $188$, the same is
true except for the first Bernstein coefficient, which factors as
\[
 b_{188,0}(y)=
 y(y+1)(y+2)^4(12y^2-21y+17).
\]
Its final quadratic is nevertheless strictly positive, since
\[
 12y^2-21y+17
 =12\left(y-\frac78\right)^2+\frac{125}{16}.
\]
Thus every Bernstein coefficient is nonnegative for $y\geq0$.  Moreover,
the last coefficient $b_{H,d}(y)$ has nonnegative monomial coefficients and
a strictly positive constant term on every row.  Therefore
$Q_H(y,\theta)>0$ for $0<\theta\leq1$: in the open interval the last
Bernstein basis function is positive, and at $\theta=1$ its coefficient is
the value of the polynomial.  Multiplication by $\theta/r^m$ proves the
claim.  At $\theta=0$ the graphon is $T_r$, giving equality.
\end{proof}

The exact degrees in \eqref{eq:interpolation-bernstein} are
\[
\begin{array}{c|l}
d&\text{Atlas IDs}\\ \hline
5&43,118,122,124,130\\
6&126,127,137,139\\
7&145,147,148,151,152,153,157,160\\
8&168,169,171,174,178\\
9&181,185,188\\
10&194,196,199\\
11&203.
\end{array}
\]
Thus the constant-versus-Turan comparison is not hiding a negative point on
their natural fixed-density segment, even at the smallest admissible Turan
grid.  This result is exact on the whole segment and does not use the
high-density estimate in Theorem~\ref{thm:regular}.

\section{Exact audit of the 29 open rows}

For the current open Atlas rows, exact edge-subset enumeration gives the
following signatures.  The last entry is the least integer $r_0$ for which
the stochastic-block corollary \eqref{eq:block-bound} is strict for every
$r\geq r_0$.
\[
\begin{array}{c|rrrrr}
H&43&118&122&124&126\\
(N_3,\Xi,r_0)&(1,55,56)&(1,126,127)&(1,123,124)&(1,125,126)&(1,112,113)\\ \hline
H&127&130&137&139&145\\
(N_3,\Xi,r_0)&(1,99,100)&(2,126,64)&(3,563,188)&(3,556,186)&(2,484,243)\\ \hline
H&147&148&151&152&153\\
(N_3,\Xi,r_0)&(2,470,236)&(2,456,229)&(1,415,416)&(2,442,222)&(2,444,223)\\ \hline
H&157&160&168&169&171\\
(N_3,\Xi,r_0)&(5,1875,376)&(5,1849,370)&(3,1609,537)&(4,1654,414)&(3,1538,513)\\ \hline
H&174&178&181&185&188\\
(N_3,\Xi,r_0)&(2,1434,718)&(7,5426,776)&(6,5022,838)&(5,4839,968)&(4,4626,1157)\\ \hline
H&194&196&199&203&\\
(N_3,\Xi,r_0)&(7,13586,1941)&(7,13377,1912)&(6,13182,2198)&(9,35258,3918)&
\end{array}
\]
Every constant is an integer obtained directly from \eqref{eq:Xi}; no
floating-point comparison is used.

In the same order, the graph6 codes are
\begin{center}
\begin{tabular}{c|c@{\qquad}c|c}
43&\verb|Dlc|&118&\verb|Eht?|\\
122&\verb|Eld?|&124&\verb|Exd?|\\
126&\verb|ERUO|&127&\verb|EZEG|\\
130&\verb|E{CW|&137&\verb|EjsG|\\
139&\verb|Eju?|&145&\verb|Exe_|\\
147&\verb|EhMg|&148&\verb|EyUG|\\
151&\verb|EhdW|&152&\verb|EhNG|\\
153&\verb|Ehf_|&157&\verb|E~`_|\\
160&\verb|E~@g|&168&\verb|E^MG|\\
169&\verb|Exf_|&171&\verb|Ehew|\\
174&\verb|EtTg|&178&\verb|En}?|\\
181&\verb|E^mG|&185&\verb|Exv_|\\
188&\verb|EzNG|&194&\verb|E~wW|\\
196&\verb|ER~g|&199&\verb|El^g|\\
203&\verb|E~^G|&&
\end{tabular}
\end{center}
Since every $r_0$ above exceeds the required $\chi(H)-1$, all displayed
densities lie in the conjecture's admissible interval.  At the limiting
endpoint $p=1$, the edge density forces $W=1$ almost everywhere, while the
monicity of $\chi_H$ gives $\Phi_H(1)=1$ by polynomial continuation; hence
$t(H,W)=\Phi_H(1)=1$ separately.

\section{Assumptions and limitations}

The proof uses only finite graphic-rank sums, leaf integration for a regular
graphon, the spectral theorem and trace identities for a self-adjoint
Hilbert--Schmidt integral operator, the elementary scalar inequality in
Lemma~\ref{lem:cycle}, and bounded convergence implicit in finite graphon
integrals.  No finite-graph approximation, optimization, regularity lemma,
or tensor search is used.

Exact degree regularity is essential to the stated spectral proof, but
no equal-cell, finite-step, or discrete-density assumption remains.  The
independent argument in
\path{notes/maximum_degree_complement_high_density_stability.tex} now covers
arbitrary nonregular complements whenever their maximum degree is uniformly
small, using degree variance and transitivity defect rather than the spectral
identity.  Complements with persistent high-degree hubs remain outside both
theorems.  The interpolation theorem removes the
high-density restriction only on the one-dimensional equal-cell segment
$B_\theta$.  The explicit general thresholds are deliberately coarse
because every higher-rank sign is discarded.  Other directions are partly
covered by the independent local and complete-multipartite theorems, but not
by this result.  No catalogue status changes.

Run
\begin{verbatim}
python codes/principled_negative_tests.py
\end{verbatim}
from the repository root.  The line headed
\begin{center}
\texttt{regular-complement graphon high-density theorem}
\end{center}
reconstructs all $29$ pairs $(N_3(H),\Xi_H)$ by exact enumeration and
compares them with the hard-coded table in the verifier.  The line headed
\texttt{full Turan-to-constant interpolation} reconstructs
\eqref{eq:interpolation-polynomial}, all Bernstein certificates, and the
Atlas $188$ quadratic sum-of-squares repair.

\end{document}
